\section{Pendahuluan}

\subsection{Latar Belakang}
\textit{Block cipher} merupakan komponen dasar kriptografi kunci simetri modern. Selain dipakai langsung
untuk menjaga kerahasiaan data, \textit{block cipher} juga menjadi bahan penyusun primitif lain, seperti
\textit{stream cipher} (melalui mode OFB dan CTR), \textit{message authentication code} (misalnya CMAC),
dan skema \textit{authenticated encryption}~\cite{rinaldi}. Karena perannya yang
mendasar tersebut, pemahaman tentang cara sebuah \textit{block cipher} dirancang, bukan hanya cara
memakainya, menjadi bekal penting dalam mempelajari kriptografi.

Landasan perancangan \textit{block cipher} modern berasal dari Shannon~\cite{rinaldi}, yang merumuskan
dua prinsip: \textit{confusion}, yaitu menyembunyikan hubungan antara kunci dan \textit{ciphertext}, dan
\textit{diffusion}, yaitu menyebarkan pengaruh setiap bit \textit{plaintext} ke seluruh
\textit{ciphertext}. Shannon juga mengusulkan \textit{product cipher}, yakni pengulangan operasi
substitusi dan transposisi yang masing-masing lemah sehingga gabungannya menjadi kuat. Hampir semua
\textit{block cipher} yang dikenal saat ini, termasuk DES dan AES, adalah \textit{iterated cipher} yang
mengikuti gagasan tersebut.

Sejarah \textit{block cipher} juga menunjukkan bahwa rancangan yang baik pada masanya dapat menjadi
tidak memadai. Kunci DES yang hanya 56 bit kini rentan terhadap \textit{exhaustive key search}
seiring kemajuan daya komputasi~\cite{rinaldi}, dan strukturnya menjadi sasaran
kriptanalisis diferensial serta kriptanalisis linier~\cite{rinaldi}. Ukuran blok 64 bit
pun terbukti bermasalah: serangan Sweet32 memanfaatkan \textit{birthday bound} $2^{32}$ blok (sekitar
32~GB data) untuk memulihkan \textit{plaintext} pada mode CBC~\cite{sweet32}. Pelajaran dari
kasus-kasus ini, yaitu kunci harus cukup panjang, blok harus cukup lebar, dan komponen nonlinier harus
tahan terhadap serangan statistik, menjadi pertimbangan langsung dalam rancangan pada laporan ini.

Tugas ini meminta perancangan sebuah \textit{block cipher} baru yang memenuhi prinsip
\textit{confusion} dan \textit{diffusion}, memakai substitusi berbasis \textit{S-box}, transposisi, serta
minimal satu operasi tambahan, dijalankan dalam beberapa putaran dengan kunci putaran yang diturunkan
dari kunci utama, dan tidak menyalin algoritma yang sudah ada. Laporan ini menyajikan \algname{},
sebuah \textit{block cipher} dengan ukuran blok 128 bit dan kunci 128 bit, beserta implementasinya pada
lima mode operasi (ECB, CBC, CFB, OFB, CTR), skema \textit{padding}, dan verifikasi integritas data.
Perlu ditekankan bahwa \algname{} dibuat untuk tujuan pembelajaran. Algoritma ini belum melalui
kriptanalisis oleh pihak lain, sehingga tidak dimaksudkan untuk menggantikan standar seperti AES pada
penggunaan nyata.

\subsection{Review Block Cipher Sejenis}
Bagian ini meninjau secara singkat beberapa \textit{block cipher} yang menjadi pembanding, terutama
dari sisi struktur, komponen nonlinier, dan mekanisme difusinya.

\paragraph{DES.} \textit{Data Encryption Standard}~\cite{rinaldi} beroperasi pada blok 64 bit dengan kunci
efektif 56 bit dan 16 putaran jaringan Feistel. Fungsi putarannya mengekspansi 32 bit
menjadi 48 bit, meng-XOR-kannya dengan kunci putaran, melewatkannya ke delapan \textit{S-box}
$6 \to 4$ bit, lalu menerapkan permutasi P. Struktur Feistel membuat enkripsi dan dekripsi memakai
rangkaian yang sama dengan urutan kunci dibalik~\cite{rinaldi}. Kelemahan utamanya adalah kunci dan blok
yang terlalu pendek untuk ukuran saat ini~\cite{rinaldi,sweet32}.

\paragraph{AES.} \textit{Advanced Encryption Standard}~\cite{rinaldi,angraini}, yang berasal dari algoritma
Rijndael, memakai blok 128 bit dengan kunci 128, 192, atau 256 bit dan 10, 12, atau 14
putaran. AES bukan Feistel, melainkan \textit{substitution-permutation network} (SPN): seluruh blok
diproses pada setiap putaran melalui SubBytes (satu \textit{S-box} $8 \times 8$ berbasis invers di
$GF(2^8)$), ShiftRows, MixColumns, dan AddRoundKey. Rancangannya mengikuti \textit{wide trail
strategy}, yang memberi batas terbukti terhadap kriptanalisis diferensial dan linier. Konsekuensinya,
setiap komponen harus memiliki invers dan dekripsi memerlukan rangkaian operasi tersendiri.

\paragraph{Blowfish.} Blowfish~\cite{siregar} adalah Feistel 16 putaran dengan blok 64 bit dan kunci
32 sampai 448 bit. Fungsi putarannya memakai empat \textit{S-box} $8 \to 32$ bit yang isinya bergantung
pada kunci, lalu menggabungkan keluarannya dengan penjumlahan modulo $2^{32}$ dan XOR. Blowfish
menunjukkan bahwa pencampuran operasi dari grup aljabar yang berbeda dapat dipakai sebagai sumber
nonlinieritas tambahan, gagasan yang sebelumnya dipakai pada IDEA. Ukuran blok 64 bit
membuatnya rentan terhadap Sweet32~\cite{sweet32}.

\paragraph{Camellia dan TEA.} Camellia menunjukkan bahwa struktur Feistel tetap layak
untuk blok 128 bit: 18 putaran untuk kunci 128 bit, dengan fungsi putaran berupa lapisan \textit{S-box}
$8 \times 8$ diikuti lapisan difusi linier. Di sisi lain, TEA adalah Feistel 64 putaran
tanpa \textit{S-box} sama sekali, yang hanya mengandalkan penjumlahan modulo $2^{32}$, pergeseran, dan
XOR, serta memakai konstanta turunan rasio emas agar setiap putaran berbeda.

Tabel~\ref{tab:perbandingan} merangkum perbandingan parameter \textit{cipher} tersebut dengan
\algname{}.

\begin{table}[H]
  \centering
  \small
  \caption{Perbandingan \algname{} dengan \textit{block cipher} sejenis}\label{tab:perbandingan}
  \setlength{\tabcolsep}{4pt}
  \begin{tabular}{>{\raggedright\arraybackslash}p{2.1cm}
                  >{\raggedright\arraybackslash}p{2.0cm}
                  >{\raggedright\arraybackslash}p{2.4cm}
                  >{\raggedright\arraybackslash}p{2.3cm}
                  >{\raggedright\arraybackslash}p{2.3cm}
                  >{\raggedright\arraybackslash}p{2.6cm}}
    \toprule
    Parameter & DES & AES-128 & Blowfish & Camellia-128 & \algname{} \\
    \midrule
    Ukuran blok & 64 bit & 128 bit & 64 bit & 128 bit & 128 bit \\
    Ukuran kunci & 56 bit & 128 bit & 32--448 bit & 128 bit & 128 bit \\
    Jumlah putaran & 16 & 10 & 16 & 18 & 16 \\
    Struktur & Feistel & SPN & Feistel & Feistel & Feistel \\
    \textit{S-box} & 8 buah, $6 \to 4$ bit, tetap
           & 1 buah, $8 \times 8$, invers $GF(2^8)$
           & 4 buah, $8 \to 32$ bit, bergantung kunci
           & 4 buah, $8 \times 8$, tetap
           & 1 buah, $8 \times 8$, tetap, hasil pencarian acak \\
    Difusi & Ekspansi E, permutasi P
           & ShiftRows, MixColumns
           & Penjumlahan mod $2^{32}$, XOR
           & Lapisan linier per \textit{byte} (XOR)
           & \textit{Carry} penjumlahan mod $2^{32}$, rotasi, permutasi bit \\
    Operasi tambahan & Ekspansi bit
           & Perkalian di $GF(2^8)$
           & Penjumlahan mod $2^{32}$
           & Fungsi $FL$/$FL^{-1}$
           & Penjumlahan dan pengurangan mod $2^{32}$, rotasi \\
    Fungsi putaran & Tanpa \textit{state}
           & Tanpa \textit{state}
           & Tanpa \textit{state}
           & Tanpa \textit{state}
           & Ber-\textit{state} (akumulator dan umpan balik) \\
    \bottomrule
  \end{tabular}
\end{table}

\subsection{Gagasan dan Pendekatan Utama}
Rancangan \algname{} bertolak dari tiga keputusan utama.

\paragraph{Blok dan kunci 128 bit.} Ukuran blok 64 bit memenuhi batas minimum spesifikasi, tetapi
\textit{birthday bound}-nya terlalu kecil untuk mode berantai seperti CBC~\cite{sweet32}. Dengan blok
128 bit, batas tersebut naik menjadi $2^{64}$ blok. Kunci 128 bit dipilih agar sepanjang ukuran blok
dan berada jauh di luar jangkauan \textit{exhaustive key search}.

\paragraph{Jaringan Feistel sebagai kerangka.} Blok dibelah menjadi dua bagian 64 bit, $L$ dan $R$,
lalu diproses sebanyak 16 putaran:
\begin{equation}
  L_{i+1} = R_i, \qquad R_{i+1} = L_i \oplus F(R_i, K_i).
\end{equation}
Pada struktur ini fungsi $F$ hanya menghasilkan nilai yang di-XOR-kan ke separuh blok lainnya, sehingga
dekripsi cukup menghitung ulang $F$ dengan urutan kunci putaran dibalik~\cite{rinaldi,sucipta}. Fungsi
$F$ tidak perlu memiliki invers. Sifat inilah yang dimanfaatkan: $F$ dapat diisi dengan proses yang
menyimpan \textit{state} dan sulit dibalik, sesuatu yang tidak mungkin dilakukan pada SPN seperti AES.
Secara teoretis, Luby dan Rackoff menunjukkan bahwa jaringan Feistel menghasilkan
permutasi pseudoacak selama fungsi putarannya pseudoacak, sehingga upaya perancangan dapat difokuskan
pada kualitas $F$.

\paragraph{Fungsi putaran berbasis modulasi sigma-delta.} Gagasan pembeda \algname{} adalah meminjam
struktur modulator sigma-delta ($\Sigma\Delta$) dari bidang pengolahan sinyal.
Modulator $\Sigma\Delta$ terdiri atas tiga bagian: integrator yang menjumlahkan selisih antara masukan
dan umpan balik, \textit{quantizer} yang memetakan isi integrator menjadi keluaran, dan jalur umpan
balik yang mengembalikan keluaran tersebut ke masukan. Ketiganya dipetakan menjadi operasi kriptografi
pada \textit{word} 32 bit:
\begin{align}
  acc      &\leftarrow acc + (x - feedback) \bmod 2^{32}, \\
  q        &\leftarrow S(acc \oplus k), \\
  feedback &\leftarrow q.
\end{align}
Integrator menjadi akumulator dengan penjumlahan modulo $2^{32}$, \textit{quantizer} menjadi
\textit{S-box} $8 \times 8$ yang diterapkan pada keempat \textit{byte} akumulator setelah dicampur
dengan kunci, dan keluaran $q$ diumpanbalikkan ke \textit{word} berikutnya. Dengan demikian satu
langkah sudah memuat tiga jenis operasi yang secara aljabar tidak saling kompatibel, yaitu penjumlahan
modulo $2^{32}$, XOR, dan substitusi tabel. \textit{S-box} menjadi sumber utama
\textit{confusion}, sedangkan \textit{carry} pada penjumlahan dan rantai umpan balik membawa pengaruh
satu \textit{word} ke \textit{word} lainnya.

Pemetaan langsung tersebut memiliki dua kelemahan yang ditangani dalam rancangan:
\begin{enumerate}
  \item Pengaruh hanya mengalir ke depan, yaitu \textit{word} kedua bergantung pada \textit{word}
        pertama tetapi tidak sebaliknya. Untuk itu $F$ menjalankan dua lintasan, maju lalu mundur,
        dengan kunci yang ditukar pada lintasan mundur sehingga setiap \textit{word} terkena kedua
        bagian kunci putaran.
  \item Bit terendah hasil penjumlahan modulo $2^{32}$ berperilaku linier (setara dengan XOR), dan
        \textit{carry} hanya merambat ke arah bit yang lebih tinggi. Untuk itu disisipkan rotasi
        bit di antara kedua lintasan, dan $F$ diakhiri dengan permutasi bit 64 bit yang menyebarkan
        kedelapan bit dari setiap \textit{byte} ke delapan \textit{byte} yang berbeda.
\end{enumerate}

Di luar fungsi putaran, dua pendekatan lain melengkapi rancangan. Pertama, \textit{S-box} tidak
diambil dari \textit{cipher} lain, melainkan dibangkitkan secara deterministik dari sejumlah kandidat
permutasi acak dan dipilih berdasarkan nonlinieritas tertinggi serta \textit{differential uniformity}
terendah, dua ukuran ketahanan terhadap kriptanalisis linier dan
diferensial. Kedua, \textit{key schedule} menurunkan 16 kunci putaran 64 bit dari kunci
utama dengan gabungan rotasi, penjumlahan konstanta turunan rasio emas yang berbeda di tiap
putaran, dan \textit{S-box} yang sama, sehingga hubungan antarkunci putaran tidak linier.

Pada tingkat aplikasi, \algname{} dijalankan pada lima mode operasi standar~\cite{rinaldi} dengan
\textit{padding} PKCS\#7 untuk mode ECB dan CBC. Untuk menjamin integritas data digunakan
CMAC yang dibangun di atas \algname{} sendiri dalam skema \textit{Encrypt-then-MAC}, dengan
kunci enkripsi dan kunci MAC yang diturunkan terpisah dari kunci utama. Rincian rancangan dibahas pada
Bab~3, sedangkan pengujian sifat \textit{avalanche} dan \textit{strict avalanche
criterion}~\cite{sucipta,angraini,prabowo} disajikan pada Bab~5.
